\documentclass[12pt,a4paper]{article}

% ── Encoding & language ──────────────────────────────────────────────────────
\usepackage[utf8]{inputenc}
\usepackage[T1]{fontenc}
\usepackage[czech]{babel}

% ── Typography ───────────────────────────────────────────────────────────────
\usepackage{lmodern}
\usepackage{microtype}

% ── Page layout ──────────────────────────────────────────────────────────────
\usepackage[margin=2.5cm, headheight=16pt]{geometry}

% ── Colors & listings ────────────────────────────────────────────────────────
\usepackage[dvipsnames]{xcolor}
\usepackage{listings}
\usepackage{lstautogobble}

\definecolor{codebg}{HTML}{F5F5F5}
\definecolor{codeframe}{HTML}{CCCCCC}
\definecolor{keyword}{HTML}{0000CC}
\definecolor{string}{HTML}{008800}
\definecolor{comment}{HTML}{888888}

\lstset{
  backgroundcolor=\color{codebg},
  frame=single,
  rulecolor=\color{codeframe},
  basicstyle=\ttfamily\small,
  keywordstyle=\color{keyword}\bfseries,
  stringstyle=\color{string},
  commentstyle=\color{comment}\itshape,
  breaklines=true,
  breakatwhitespace=true,
  autogobble=true,
  showstringspaces=false,
  numbers=left,
  numberstyle=\tiny\color{gray},
  numbersep=6pt,
  xleftmargin=14pt,
  literate={├}{{\textbar}}1 {─}{-}1 {└}{{\texttt{+}}}1 {──}{{--}}2
}

% ── Rename listing caption ────────────────────────────────────────────────────
\renewcommand{\lstlistingname}{Ukázka}

% ── Hyperlinks ────────────────────────────────────────────────────────────────
\usepackage[colorlinks=true, linkcolor=MidnightBlue, urlcolor=MidnightBlue]{hyperref}

% ── Misc ──────────────────────────────────────────────────────────────────────
\usepackage{booktabs}
\usepackage{enumitem}
\usepackage{parskip}
\usepackage{titlesec}
\usepackage{fancyhdr}
\usepackage{graphicx}
\usepackage{tcolorbox}
\tcbuselibrary{skins,breakable}

% ── Callout boxes ─────────────────────────────────────────────────────────────
\newtcolorbox{infobox}{
  colback=blue!5!white, colframe=MidnightBlue,
  fonttitle=\bfseries, title=Poznámka,
  breakable
}
\newtcolorbox{warnbox}{
  colback=orange!8!white, colframe=BurntOrange,
  fonttitle=\bfseries, title=Upozornění,
  breakable
}

% ── Header / footer ───────────────────────────────────────────────────────────
\pagestyle{fancy}
\fancyhf{}
\rhead{Hub -- uživatelský manuál}
\lhead{\leftmark}
\cfoot{\thepage}
\renewcommand{\headrulewidth}{0.4pt}

% ── Title ─────────────────────────────────────────────────────────────────────
\title{
  \Huge\bfseries Hub\\[6pt]
  \large Uživatelský manuál\\[4pt]
  \normalsize Backend systému pro spouštění webových aplikací (her)
}
\author{}
\date{\today}

% ─────────────────────────────────────────────────────────────────────────────
\begin{document}
\maketitle
\thispagestyle{empty}
\vspace{1cm}

\clearpage

\tableofcontents
\clearpage

% ═══════════════════════════════════════════════════════════════════════════════
\clearpage
\section{Přehled systému}
% ═══════════════════════════════════════════════════════════════════════════════

Hub je centrální backendová služba, která zajišťuje přístup klientů k webovým aplikacím (hrám) prostřednictvím \textbf{reverse proxy}. Celý systém se skládá ze tří vrstev:

\begin{enumerate}[label=\arabic*.]
  \item \textbf{Backend (hub-backend)} -- Go aplikace; obsluhuje REST API pro registraci her a provádí reverse proxy na příslušnou herní službu.
  \item \textbf{Frontend (hub-frontend)} -- Webové rozhraní pro uživatele; zobrazuje seznam dostupných her.
  \item \textbf{Herní kontejnery} -- Libovolné aplikace/hry běžící v Docker kontejnerech, připojené ke sdílené síti \texttt{hub-shared}.
\end{enumerate}

\subsection{Tok požadavku}

\begin{enumerate}[label=\arabic*.]
  \item Prohlížeč odešle požadavek na adresu \texttt{https://\textit{DOMÉNA}/games/\textit{id}/...}
  \item Backend najde registrovanou hru podle \texttt{id}.
  \item Požadavek je přeposlán na interní adresu herního kontejneru (např.\ \texttt{http://mygame:8080}).
  \item Odpověď hry je vrácena prohlížeči.
\end{enumerate}

Požadavky na ostatní cesty (vše mimo \texttt{/games/} a \texttt{/api/}) jsou přeposílány na frontend službu.

\subsection{REST API -- přehled endpointů}

\begin{center}
\begin{tabular}{@{}lll@{}}
\toprule
\textbf{Metoda} & \textbf{Cesta} & \textbf{Popis} \\
\midrule
GET    & \texttt{/api/games}        & Vrátí seznam registrovaných her \\
POST   & \texttt{/api/games}        & Zaregistruje novou hru \\
DELETE & \texttt{/api/games/\{id\}} & Odregistruje hru \\
ANY    & \texttt{/games/\{id\}/...} & Reverse proxy na danou hru \\
ANY    & \texttt{/}                 & Reverse proxy na frontend \\
\bottomrule
\end{tabular}
\end{center}

% ═══════════════════════════════════════════════════════════════════════════════
\clearpage
\section{Nasazení systému Hub na server}
% ═══════════════════════════════════════════════════════════════════════════════

\subsection{Požadavky}

\begin{itemize}
  \item Docker Engine \(\geq\) 24 a Docker Compose Plugin (\texttt{docker compose})
  \item Git
  \item Doménové jméno směrující na server (pro produkci s HTTPS)
  \item Otevřené porty 80 a 443
\end{itemize}

\subsection{Stažení repozitářů}

Repozitáře backendu a frontendu \textbf{musí ležet vedle sebe} ve stejném nadřazeném adresáři, protože \texttt{docker-compose.yml} backendu odkazuje na frontend cestu \texttt{../hub-frontend}.

\begin{lstlisting}[language=bash, caption={Klonovani repozitaru}]
  mkdir hub && cd hub
  git clone https://github.com/<org>/hub-backend.git   backend
  git clone https://github.com/<org>/hub-frontend.git  hub-frontend
\end{lstlisting}

Výsledná struktura adresářů:

\begin{lstlisting}[numbers=none]
  hub/
  ├── backend/          # repozitar backendu
  └── hub-frontend/     # repozitar frontendu
\end{lstlisting}

\subsection{Konfigurace prostředí}

Zkopírujte vzorový soubor a upravte hodnoty:

\begin{lstlisting}[language=bash]
  cd backend
  cp .env.example .env
  nano .env           # nebo vim, code, ...
\end{lstlisting}

Klíčové proměnné:

\begin{center}
\begin{tabular}{@{}lll@{}}
\toprule
\textbf{Proměnná} & \textbf{Výchozí hodnota} & \textbf{Popis} \\
\midrule
\texttt{APP\_ENV}          & \texttt{development}           & Režim běhu (\texttt{production} / \texttt{development}) \\
\texttt{DOMAIN}            & --                             & FQDN domény (např.\ \texttt{hub.example.com}) \\
\texttt{CERT\_EMAIL}       & --                             & E-mail pro Let's Encrypt \\
\texttt{CERT\_CACHE\_DIR}  & \texttt{/certs}                & Adresář pro cache TLS certifikátů \\
\texttt{FRONTEND\_TARGET}  & \texttt{http://frontend:5173}  & Interní URL frontend kontejneru \\
\texttt{HTTP\_ADDR}        & \texttt{:80}                   & HTTP port (redirect + ACME challenge) \\
\texttt{HTTPS\_ADDR}       & \texttt{:443}                  & HTTPS port \\
\bottomrule
\end{tabular}
\end{center}

Příklad \texttt{.env} pro produkci:

\begin{lstlisting}[caption={Příklad .env pro produkci}]
  APP_ENV=production
  DOMAIN=hub.example.com
  CERT_EMAIL=admin@example.com
  CERT_CACHE_DIR=/certs
  FRONTEND_TARGET=http://frontend:5173
  HTTP_ADDR=:80
  HTTPS_ADDR=:443
\end{lstlisting}

\subsection{Vytvoření sdílené Docker sítě}

Sdílená síť \texttt{hub-shared} musí existovat \textbf{před} spuštěním compose, protože je deklarována jako \texttt{external: true}.

\begin{lstlisting}[language=bash]
  docker network create hub-shared
\end{lstlisting}

Tuto síť stačí vytvořit jednou; po restartu serveru ji Docker zachová.

\subsection{Spuštění systému}

\begin{lstlisting}[language=bash, caption={Sestaveni a spusteni hub (backend + frontend)}]
  cd backend
  docker compose up --build -d
\end{lstlisting}

Příznak \texttt{--build} zajistí sestavení obrazů ze zdrojového kódu. Příznak \texttt{-d} spustí služby na pozadí.

\begin{infobox}
Při prvním spuštění v produkčním režimu backend automaticky získá TLS certifikát od Let's Encrypt. Port 80 musí být dostupný z internetu kvůli ACME HTTP-01 challenge.
\end{infobox}

\subsection{Instalační skript}

Pro zjednodušení nasazení je k dispozici skript \texttt{install-hub.sh} (viz příloha nebo kořen repozitáře). Skript:

\begin{enumerate}[label=\arabic*.]
  \item Naklonuje oba repozitáře do správné struktury.
  \item Vytvoří Docker síť \texttt{hub-shared} (pokud neexistuje).
  \item Zkopíruje \texttt{.env.example} $\rightarrow$ \texttt{.env} (pokud \texttt{.env} neexistuje).
  \item Spustí \texttt{docker compose up --build -d}.
\end{enumerate}

Použití:

\begin{lstlisting}[language=bash]
  bash install-hub.sh
  # nebo
  chmod +x install-hub.sh && ./install-hub.sh
\end{lstlisting}

\begin{warnbox}
Před spuštěním skriptu upravte \texttt{.env} s hodnotami pro vaše prostředí (zejména \texttt{DOMAIN} a \texttt{CERT\_EMAIL}).
\end{warnbox}

% ═══════════════════════════════════════════════════════════════════════════════
\clearpage
\section{Integrace aplikace (hry) do systému Hub}
% ═══════════════════════════════════════════════════════════════════════════════

Tato kapitola popisuje, jak vytvořit novou hru/aplikaci kompatibilní s Hub systémem. Pro úspěšnou integraci musí aplikace splňovat tři podmínky:

\begin{enumerate}[label=\arabic*.]
  \item Běžet v Docker kontejneru.
  \item Být připojena na sdílenou síť \texttt{hub-shared}.
  \item Registrovat se u backendu přes REST API při startu.
\end{enumerate}

\subsection{Požadavek 1: Práce za reverse proxy}
\label{sec:revproxy}

Backend přeposílá požadavky na hru a \textbf{ořezává} prefix \texttt{/games/\{id\}} z URL. Aplikace tedy vidí cestu \textbf{od kořene} (\texttt{/}).

\begin{infobox}
Příklad: uživatel navštíví \texttt{https://hub.example.com/games/chess/board} -- hra obdrží požadavek na cestu \texttt{/board}.
\end{infobox}

Důsledky pro aplikaci:

\begin{itemize}
  \item Všechny interní odkazy (HTML \texttt{href}, JS \texttt{fetch}, CSS \texttt{url()}) musí být \textbf{relativní} nebo začínat \texttt{/games/\{id\}/}.
  \item Statické soubory (JS, CSS, obrázky) musí být servovány ze stejné cesty jako HTML -- ideálně z kořene.
  \item Pokud aplikace používá klientský router (SPA), musí server vracet \texttt{index.html} pro všechny neznámé cesty (catch-all).
\end{itemize}

Příklad konfigurace Nginx pro SPA aplikaci:

\begin{lstlisting}[language=bash, caption={nginx.conf -- catch-all pro SPA}]
  server {
      listen 80;
      root /usr/share/nginx/html;
      index index.html;

      location / {
          try_files $uri $uri/ /index.html;
      }
  }
\end{lstlisting}

\subsection{Požadavek 2: Docker kontejner a sdílená síť}

Hra musí být kontejnerizována a připojena do sítě \texttt{hub-shared}. Hub i hra pak mohou komunikovat přes Docker DNS (jménem kontejneru / služby).

Minimální \texttt{docker-compose.yml} pro hru:

\begin{lstlisting}[language=bash, caption={docker-compose.yml herní aplikace}]
  services:
    mygame:
      build: .
      expose:
        - "80"        # port exponovan jen uvnitr site, NE na hostiteli
      networks:
        - hub-shared
      restart: unless-stopped

  networks:
    hub-shared:
      external: true  # pripojeni na existujici sdilenou sit hubu
\end{lstlisting}

\begin{warnbox}
Používejte \texttt{expose} (interní port), nikoli \texttt{ports} (publikovaný port na hostiteli). Hra nemá být dostupná přímo z internetu -- veškerý přístup musí procházet přes Hub.
\end{warnbox}

\subsection{Požadavek 3: Registrace u backendu}

Při startu kontejneru se hra musí zaregistrovat u Hub backendu pomocí HTTP POST na \texttt{/api/games}. Tuto registraci provádí startovací skript.

\subsubsection{Registrační skript}
\label{sec:regscript}

Vytvořte soubor \texttt{register-and-run.sh}:

\begin{lstlisting}[language=bash, caption={register-and-run.sh}]
  #!/bin/sh
  set -eu

  BACKEND_URL=${GAME_HUB_URL:-http://backend:8080}
  ENDPOINT="$BACKEND_URL/api/games"

  # Upravte pole dle sve aplikace
  payload='{
    "id":          "mygame",
    "name":        "My Game",
    "target":      "http://mygame:80",
    "status":      "online",
    "icon":        "nejaka emoji",
    "author":      "Jmeno autora",
    "description": "Strucny popis hry.",
    "image":       "https://example.com/mygame-icon.png"
  }'

  max_retries=30
  sleep_seconds=2
  attempt=1

  while [ $attempt -le $max_retries ]; do
    echo "[register] Pokus $attempt: POST $ENDPOINT"
    http_code=$(curl -s -o /tmp/resp -w "%{http_code}" \
      -X POST "$ENDPOINT" \
      -H "Content-Type: application/json" \
      -d "$payload" || true)

    if [ "$http_code" = "201" ] || [ "$http_code" = "200" ]; then
      echo "[register] Registrace uspela (HTTP $http_code)"
      break
    fi

    echo "[register] Selhani (HTTP $http_code), dalsi pokus za ${sleep_seconds}s"
    attempt=$((attempt + 1))
    sleep $sleep_seconds
  done

  # Spustit vlastni aplikaci (nginx, node, ...)
  exec nginx -g "daemon off;"
\end{lstlisting}

\begin{infobox}
Pole \texttt{target} musí obsahovat interní Docker adresu hry -- tedy název služby a port, pod kterým ji Hub backend dosáhne uvnitř sítě \texttt{hub-shared}. Například \texttt{http://mygame:80}.
\end{infobox}

\subsubsection{Struktura JSON záznamu hry}

\begin{center}
\begin{tabular}{@{}lll@{}}
\toprule
\textbf{Pole} & \textbf{Typ} & \textbf{Popis} \\
\midrule
\texttt{id}          & string & Unikátní identifikátor (URL-safe, bez mezer) \\
\texttt{name}        & string & Zobrazovaný název hry \\
\texttt{target}      & string & Interní URL kontejneru hry (Docker DNS) \\
\texttt{status}      & string & Stav hry (\texttt{online} / \texttt{offline}) \\
\texttt{icon}        & string & Emoji nebo krátký symbol \\
\texttt{author}      & string & Jméno autora \\
\texttt{description} & string & Krátký popis \\
\texttt{image}       & string & URL náhledového obrázku (volitelné) \\
\bottomrule
\end{tabular}
\end{center}

\clearpage

\subsubsection{Integrace skriptu do Dockerfile}
\label{sec:dockerfile}

\begin{lstlisting}[language=bash, caption={Dockerfile herní aplikace (zkrácený příklad)}]
  FROM nginx:alpine

  # curl je potreba pro registracni skript
  RUN apk add --no-cache curl

  # Zkopirujte vysledny web do Nginx rootu
  COPY dist/index.html /usr/share/nginx/html/index.html
  COPY nginx.conf /etc/nginx/conf.d/default.conf

  # Registracni skript
  COPY register-and-run.sh /register-and-run.sh
  RUN chmod +x /register-and-run.sh

  EXPOSE 80

  # Skript nejprve zaregistruje hru, pak spusti nginx
  CMD ["/register-and-run.sh"]
\end{lstlisting}

\subsection{Přehled proměnné prostředí pro herní kontejner}

\begin{center}
\begin{tabular}{@{}lll@{}}
\toprule
\textbf{Proměnná} & \textbf{Výchozí hodnota} & \textbf{Popis} \\
\midrule
\texttt{GAME\_HUB\_URL} & \texttt{http://backend:8080} & URL Hub backendu uvnitř Docker sítě \\
\bottomrule
\end{tabular}
\end{center}

% ═══════════════════════════════════════════════════════════════════════════════
\clearpage
\section{Příklad -- kompletní integrace hry}
% ═══════════════════════════════════════════════════════════════════════════════

Tato sekce shrnuje celý postup na ukázkovém příkladu hry \textit{Chess}.

\begin{enumerate}[label=\textbf{Krok \arabic*.},leftmargin=*]

  \item \textbf{Ujistěte se, že síť \texttt{hub-shared} existuje:}
\begin{lstlisting}[language=bash, numbers=none]
  docker network create hub-shared
\end{lstlisting}

  \item \textbf{Vytvořte registrační skript} \texttt{register-and-run.sh} (viz \hyperref[sec:regscript]{sekci~\ref*{sec:regscript}}) s hodnotami pro vaši hru (\texttt{id=chess}, \texttt{target=http://chess:80}, \ldots).

  \item \textbf{Nakonfigurujte Nginx} (\texttt{nginx.conf}) s catch-all pro SPA (viz \hyperref[sec:revproxy]{sekci~\ref*{sec:revproxy}}).

  \item \textbf{Napište Dockerfile} zahrnující \texttt{curl}, registrační skript a Nginx (viz \hyperref[sec:dockerfile]{sekci~\ref*{sec:dockerfile}}).

  \item \textbf{Vytvořte \texttt{docker-compose.yml}} hry připojené na \texttt{hub-shared}:
\begin{lstlisting}[language=bash, numbers=none]
  services:
    chess:
      build: .
      expose:
        - "80"
      networks:
        - hub-shared
      restart: unless-stopped

  networks:
    hub-shared:
      external: true
\end{lstlisting}

  \item \textbf{Spusťte hru:}
\begin{lstlisting}[language=bash, numbers=none]
  docker compose up --build -d
\end{lstlisting}

  \item \textbf{Ověřte registraci:}
\begin{lstlisting}[language=bash, numbers=none]
  curl https://hub.example.com/api/games
\end{lstlisting}
  V odpovědi by měla být hra \texttt{chess} se stavem \texttt{online}.

  \item \textbf{Hra je dostupná} na adrese \texttt{https://hub.example.com/games/chess}.

\end{enumerate}

% ═══════════════════════════════════════════════════════════════════════════════
\clearpage
\section{Odregistrování hry}
% ═══════════════════════════════════════════════════════════════════════════════

Pokud chcete hru z Hubu odebrat (např.\ při odstavení kontejneru), zavolejte DELETE endpoint:

\begin{lstlisting}[language=bash]
  curl -X DELETE https://hub.example.com/api/games/chess
\end{lstlisting}

Odpověď: \texttt{204 No Content}. Hra zmizí ze seznamu dostupných her.

\begin{infobox}
Registrace her je ukládána do souboru \texttt{data/games.json} uvnitř backend kontejneru. Po restartu backendu se dříve registrované hry automaticky obnoví. Herní kontejnery se musí po restartu samy znovu zaregistrovat (registrační skript toto zajišťuje automaticky při spuštění kontejneru).
\end{infobox}

% ═══════════════════════════════════════════════════════════════════════════════
\clearpage
\section{Řešení problémů}
% ═══════════════════════════════════════════════════════════════════════════════

\begin{center}
\begin{tabular}{@{}p{5cm}p{9cm}@{}}
\toprule
\textbf{Problém} & \textbf{Řešení} \\
\midrule
Hra se neregistruje & Zkontrolujte, zda je kontejner hry ve stejné síti jako backend (\texttt{hub-shared}). Ověřte DNS: \texttt{docker exec <kontejner> ping backend}. \\
\addlinespace
404 na \texttt{/games/id} & Hra není registrována. Zkontrolujte logy kontejneru hry a výstup registračního skriptu. \\
\addlinespace
Aktiva (CSS/JS) se nenačítají & Statické soubory musí být na relativních cestách. Ujistěte se, že Nginx vrací \texttt{index.html} pro neznámé cesty (catch-all). \\
\addlinespace
Let's Encrypt selhává & Port 80 musí být dostupný z internetu. Zkontrolujte firewall a zda \texttt{DOMAIN} správně ukazuje na server. \\
\addlinespace
Síť \texttt{hub-shared} neexistuje & Spusťte \texttt{docker network create hub-shared} před \texttt{docker compose up}. \\
\bottomrule
\end{tabular}
\end{center}

\end{document}
